%======================================================================
\chapter{Correspondence with the formal proof}\label{app:formal}
%======================================================================

The table lists the definitions and results of this volume that carry a ``Formal counterpart''
line, together with the corresponding declarations of the Lean development. All declaration
names start with \lean{HPC.}; file names are relative to the root of the repository. The formal
statements are sometimes slightly more general, or are split differently into lemmas; the table
points to the place where the content of each result is proved. Rows are in the order of the
text.

The main theorems of the volume are, in the Lean development (prefix \lean{HPC.} omitted):
\begin{center}
\begin{tabular}{@{}p{0.42\textwidth}p{0.54\textwidth}@{}}
\toprule
Uniqueness, with two pairs of points (Theorem~\ref{thm:unique-points}) & \lean{Plane.theorem1_HIC}\\
Uniqueness (Theorem~\ref{thm:unique}) & \lean{Plane.theorem1_HIC'}\\
Homogeneity (Theorem~\ref{thm:homog}) & \lean{Plane.homogeneous_of_isPseudoArc}\\
Every pseudo-arc is the plane one (Corollary~\ref{cor:plane-unique}) & \lean{Plane.pseudoArc_unique}\\
Subcontinua (Corollary~\ref{cor:sub-homeo}) & \lean{Plane.homeomorphic_subcontinuum}\\
Two descriptions (Theorem~\ref{thm:two-desc}) & \lean{theorem12}\\
Every pseudo-arc has a description (Theorem~\ref{thm:desc-exists}) &
\lean{Plane.exists_description_of_chainable_HI}\\
\bottomrule
\end{tabular}
\end{center}
All of them are proved in the Lean development without \lean{sorry} and without axioms beyond the
standard ones of Lean and Mathlib.

\begin{longtable}{@{}p{0.40\textwidth}p{0.56\textwidth}@{}}
\toprule
In these notes & In the Lean development\\
\midrule
\endhead
\bottomrule
\endfoot
\multicolumn{2}{@{}l}{\emph{Chapter~\ref{ch:composants}}}\\
Definition~\ref{def:composant} & \lean{HPC.composant} (in \lean{HPC/PseudoArc/Statements79.lean}).\\
Lemma~\ref{lem:ne-of-composant} & \lean{HPC.Plane.PA51.ne_of_not_mem_composant}.\\
Boundary bumping, clopen form (Lemma~\ref{lem:bumping}) & \lean{HPC.T13.clopen_meets_closure_diff}.\\
Lemma~\ref{lem:union-clopen} & \lean{HPC.T13.isPreconnected_union_clopen}.\\
Theorem~\ref{thm:thin} & \lean{HPC.T13.false_of_interior}.\\
Theorem~\ref{thm:composant-proper} & \lean{HPC.T13.exists_not_mem_proper_subcontinua} (in a compact ambient space) and \lean{HPC.composant_ne_of_indecomposable}.\\
Corollary~\ref{cor:other-composant} & \lean{HPC.Plane.exists_not_mem_composant}. Formally the theorem is applied in the compact space $M$, and the subcontinuum found there is transported back to $X$.\\
Lemma~\ref{lem:sub-HI} & \lean{HPC.Plane.PA51.hI_sub}, \lean{HPC.Plane.PA51.not_mem_composant_sub}.\\
Lemma~\ref{lem:cover-meeting} & (a) is \lean{HPC.Plane.exists_cover_meeting}; (b) is \lean{HPC.Plane.PA51.inter_inter_nonempty} and \lean{HPC.Plane.PA51.subChain}.\\
Corollary~\ref{cor:sub-snake} & \lean{HPC.Plane.PA51.snake_sub}.\\[4pt]
\multicolumn{2}{@{}l}{\emph{Chapter~\ref{ch:consolidations}}}\\
Definition~\ref{def:consolidation} & \lean{HPC.Chain.IsConsolidationOf}.\\
Lemma~\ref{lem:block-chain} & \lean{HPC.Chain.IsBlockMap}, \lean{HPC.Chain.ofBlocks}, \lean{HPC.Chain.blockLink_adjacent}, \lean{HPC.Chain.ofBlocks_isConsolidationOf}.\\
Lemma~\ref{lem:block-from} & \lean{HPC.Fold.Q78.isFrom_ofBlocks_iff}.\\
Lemma~\ref{lem:schedule} & \lean{HPC.IsSchedule}, \lean{HPC.IsSchedule.exists_blockMap}.\\
Lemma~\ref{lem:consolidate-length} & \lean{HPC.Chain.exists_consolidation_isFrom}.\\
Lemma~\ref{lem:index-dist} & \lean{HPC.Chain.index_dist_le_of_subset}.\\
Bing 1948, Theorem~1 (Theorem~\ref{thm:bing1}) & \lean{HPC.Chain.isCrookedIn'_of_isConsolidationOf}; the interpolation of Step~4 is \lean{HPC.Chain.exists_value_between}.\\
Lemma~\ref{lem:reverse} & \lean{HPC.Chain.reverse}, \lean{HPC.Chain.IsCrookedIn'.reverse_left}, \lean{HPC.Chain.IsCrookedIn'.reverse_right}.\\
Lemma~\ref{lem:minidx} & \lean{HPC.Chain.Refines.minIdx}, \lean{HPC.Chain.Refines.maxIdx}, \lean{HPC.Chain.subset_adjacent_of_subset}, \lean{HPC.Chain.minIdx_lip}, \lean{HPC.Chain.maxIdx_lip}, \lean{HPC.Chain.minIdx_near_of_inter}.\\
Lemma~\ref{lem:span-chained} & \lean{HPC.Chain.span}, \lean{HPC.Chain.chainedAt_span}.\\
Definition~\ref{def:subchain-from} & \lean{HPC.Chain.IsSubchainFrom}.\\
Lemma~\ref{lem:subchain-exists} & \lean{HPC.Chain.exists_isSubchainFrom}.\\[4pt]
\multicolumn{2}{@{}l}{\emph{Chapter~\ref{ch:patterns}}}\\
Pattern following; Bing 1948, Theorem~6 (Theorem~\ref{thm:pattern-following}) & \lean{HPC.Description.patternFollowing}; the property itself is \lean{HPC.Description.PatternFollowing}, and patterns are \lean{HPC.Description.IsPattern}; the proof is \lean{HPC.Description.theorem6Pos_of_ne}, with \lean{HPC.Description.patternFollowing_of_theorem6Pos} for the passage to arbitrary $i$ (via the shifted description \lean{HPC.Description.shift}).\\
Remark~\ref{rem:pf-false} & \lean{HPC.Description.unitDescription}, \lean{HPC.Description.not_patternFollowing_unconditional}.\\
Lemma~\ref{lem:allowed} & \lean{HPC.PatCore.Setting}, \lean{HPC.PatCore.Slow}, \lean{HPC.PatCore.Al}, \lean{HPC.PatCore.al_canon}, \lean{HPC.PatCore.al_move}, \lean{HPC.PatCore.al_near}, \lean{HPC.PatCore.al_far}, \lean{HPC.PatCore.time_ge}; \lean{HPC.PatCore.ivt_between}.\\
Definition~\ref{def:lift} & \lean{HPC.PatCore.Real}.\\
Operations on lifts (Lemma~\ref{lem:lift-ops}) & \lean{HPC.PatCore.real_glue}, \lean{HPC.PatCore.real_widen}, \lean{HPC.PatCore.real_of_neg}, \lean{HPC.PatCore.real_reflect}, \lean{HPC.PatCore.real_step}, \lean{HPC.PatCore.real_ramp}.\\
Definition~\ref{def:good} & \lean{HPC.PatCore.Good}.\\
Main induction (Theorem~\ref{thm:pf-main}) & \lean{HPC.PatCore.main_step} with the cases \lean{HPC.PatCore.case_const}, \lean{HPC.PatCore.case_plateau}, \lean{HPC.PatCore.case_plateau_end}, \lean{HPC.PatCore.case_amp1}, \lean{HPC.PatCore.case_return}, \lean{HPC.PatCore.case_noreturn}; the induction is \lean{HPC.PatCore.ih_all}.\\
Combinatorial core of pattern following (Theorem~\ref{thm:realize}) & \lean{HPC.PatCore.realize}.\\
Lemma~\ref{lem:n-tendsto} & \lean{HPC.Description.n_tendsto_atTop} (via the more general \lean{HPC.Chain.diam_le_length_mul}).\\
Thick readings (Lemma~\ref{lem:thick}) & \lean{HPC.Description.eventually_thick}.\\
Lemma~\ref{lem:setting-of-crooked} & \lean{HPC.Chain.setting_of_crooked}.\\
Lemma~\ref{lem:consolidation-of-lift} & \lean{HPC.Chain.exists_consolidation_of_lift}.\\[4pt]
\multicolumn{2}{@{}l}{\emph{Chapter~\ref{ch:matching}}}\\
Definition~\ref{def:matching} & \lean{HPC.MatchingCovers}, \leansub{HPC.MatchingCovers.Tracks}{1}, \leansub{HPC.MatchingCovers.Tracks}{2}.\\
Matching covers; Bing 1948, Theorem~11 (Theorem~\ref{thm:matching}) & \lean{HPC.MatchingCovers.exists_homeomorph}.\\
Fattening & \lean{HPC.tail}, \lean{HPC.rad}, \lean{HPC.mesh}, \lean{HPC.fat}, \leansub{HPC.MatchingCovers.y}{1}, \leansub{HPC.MatchingCovers.y}{2}, \lean{HPC.diam_fat_le_mesh}, \lean{HPC.mesh_tendsto_zero}.\\
Nesting (Lemma~\ref{lem:nesting}) & \lean{HPC.closure_fat_subset_fat}.\\
Step 1 (Lemma~\ref{lem:step1}) & \leansub{HPC.MatchingCovers.step1}{2}.\\
Step 2 (Lemma~\ref{lem:step2}) & \leansub{HPC.MatchingCovers.step2}{2}.\\
Step 3 (Lemma~\ref{lem:limit-point}) & \lean{HPC.MatchingCovers.exists_limit_point}.\\
Step 4: independence of choices (Lemma~\ref{lem:T-indep}) & \lean{HPC.MatchingCovers.toFun}, \leansub{HPC.MatchingCovers.toFun_mem_y}{2}.\\
Step 5 (Lemma~\ref{lem:T-cont}) & \lean{HPC.MatchingCovers.continuous_toMap}.\\
Step 6 (Lemma~\ref{lem:T-inj}) & \lean{HPC.MatchingCovers.injective_toMap}.\\
Step 7 (Lemma~\ref{lem:T-surj}) & \lean{HPC.MatchingCovers.surjective_toMap}.\\[4pt]
\multicolumn{2}{@{}l}{\emph{Chapter~\ref{ch:two}}}\\
Bing 1948, Theorem~12 (Theorem~\ref{thm:two-desc}) & \lean{HPC.theorem12}, which combines \lean{HPC.Description.exists_homeomorph} (the theorem assuming pattern following) with \lean{HPC.Description.patternFollowing}.\\
Lemma~\ref{lem:pattern-of-refines} & \lean{HPC.Chain.exists_pattern_of_refines}.\\
Matching (Lemma~\ref{lem:pattern-matching}) & \lean{HPC.Chain.pattern_matching}.\\
Lemma~\ref{lem:consolidation-meets} & \lean{HPC.Description.M_subset_consolidation}, \lean{HPC.Description.consolidation_link_inter_M_nonempty}, \lean{HPC.Description.link_inter_M_nonempty}.\\
Lemma~\ref{lem:level-refining} & \lean{HPC.Description.exists_level_refining}.\\
Definition~\ref{def:rung} & \lean{HPC.Rung}. Formally, (R3) asks that the active side be able to follow patterns (\lean{HPC.Description.CanFollow}), which holds for every level by pattern following (\lean{HPC.Description.canFollow_level}).\\
The first rung (Lemma~\ref{lem:base-rung}) & \lean{HPC.Rung.exists_base}.\\
The next rung (Lemma~\ref{lem:next-rung}) & \lean{HPC.Rung.exists_next}, via \lean{HPC.Description.exists_half_step}.\\
Corollary~\ref{cor:two-desc-K} & \lean{HPC.Plane.theorem12_of}, using the restriction \lean{HPC.Description.restrict} (in \lean{HPC/Plane/Transfer.lean}).\\[4pt]
\multicolumn{2}{@{}l}{\emph{Chapter~\ref{ch:folding}}}\\
Folding theorem; Bing 1948, Theorem~4 (Theorem~\ref{thm:folding}) & \lean{HPC.Fold.theorem4_pair} (in \lean{HPC/Fold/Theorem4Pair.lean}). The version with single links $r$ and $s$ in place of the pairs is \lean{HPC.Fold.theorem4}. In the proof, the main case is \lean{HPC.Fold.B79.theorem4_pair_core}, and the reflection uses \lean{HPC.Fold.B78.unionRange_reverse}.\\
Corollary~\ref{cor:block-chain} & \lean{HPC.Chain.exists_blockChain} (in \lean{HPC/PseudoArc/T13Step.lean}).\\
Lemma~\ref{lem:readings} & \lean{HPC.Fold.IsCrookedIvF} and \lean{HPC.Fold.isCrookedIvF_minIdx_maxIdx} package these facts with crookedness; see also \lean{HPC.Fold.B79.seq_facts}.\\
Crooked sequence (Definition~\ref{def:crooked-seq}) & \lean{HPC.Fold.Z77.CrkOn}.\\
Lemma~\ref{lem:crooked-seq} & \lean{HPC.Fold.Z77.crkOn_of_isCrookedIvF}, \lean{HPC.Fold.Z77.CrkOn.reverse}, \lean{HPC.Fold.B79.minIdx_reverse}.\\
Fitting (Definition~\ref{def:fit}) & \lean{HPC.Fold.Z77.Env}, \lean{HPC.Fold.Z77.arr}, \lean{HPC.Fold.Z77.Fit}.\\
Lemma~\ref{lem:legpos} & \lean{HPC.Fold.Z77.jwalk}, \lean{HPC.Fold.Z77.legPos}, \lean{HPC.Fold.Z77.jwalk_legPos}, \lean{HPC.Fold.Z77.legPos_lip}, \lean{HPC.Fold.Z77.legPos_eq_zero}.\\
Folds from fitting (Lemma~\ref{lem:fold-of-fit}) & \lean{HPC.Fold.Z77.fold_of_fit}, \lean{HPC.Fold.Z77.fold_of_fit_zero}.\\
Up phase (Lemma~\ref{lem:up-phase}) & \lean{HPC.Fold.Z77.up_step} (with \lean{HPC.Fold.Z77.transfer} and \lean{HPC.Fold.Z77.exists_armVal}).\\
Down phase (Lemma~\ref{lem:down-phase}) & \lean{HPC.Fold.Z77.down_step}.\\
Up core (Lemma~\ref{lem:up-core}) & \lean{HPC.Fold.Z77.up_core}.\\
Down core (Lemma~\ref{lem:down-core}) & \lean{HPC.Fold.Z77.down_core}.\\
Up states are not stuck (Lemma~\ref{lem:nofail-up}) & \lean{HPC.Fold.Z77.nofail_up}.\\
Down states after an up phase are not stuck (Lemma~\ref{lem:nofail-down}) & \lean{HPC.Fold.Z77.nofail_down}.\\
First direction (Lemma~\ref{lem:first-direction}) & \lean{HPC.Fold.Z77.first_direction}; the two conditions are \lean{HPC.Fold.Z77.OKBelow} and \lean{HPC.Fold.Z77.OKAbove}.\\
Main induction (Proposition~\ref{prop:main-up}) & \lean{HPC.Fold.Z77.main_up}.\\
Zigzag theorem (Theorem~\ref{thm:zigzag}) & \lean{HPC.Fold.Z77.greedy_fit_below}, \lean{HPC.Fold.Z77.greedy_fit_above}, and \lean{HPC.Fold.Z77.greedy_fit} (which uses Lemma~\ref{lem:first-direction} to choose $\epsilon$).\\
Bing's Theorem~3, three-link form (Corollary~\ref{cor:theorem3}) & \lean{HPC.Fold.theorem3}, via \lean{HPC.Fold.Z77.s3_zig} and \lean{HPC.exists_blocks_of_walk}.\\
Folds with a fake end link (Section~\ref{sec:fold-fake}) & \lean{HPC.Fold.B79.EvUp} and \lean{HPC.Fold.B79.EvDown}; reflecting $f$ exchanges them (\lean{HPC.Fold.B79.evUp_reflect}).\\
Crooked extensions (Lemma~\ref{lem:crooked-ext}) & \lean{HPC.Fold.B79.crkOn_lowExt}, \lean{HPC.Fold.B79.crkOn_highExt} (on sequences re-indexed so that the base is at $0$).\\
Fold with a fake low end (Lemma~\ref{lem:fake-low}) & \lean{HPC.Fold.B79.finfold_low}, from \lean{HPC.Fold.B79.fold_low_end_mode} (the version on $\Z$-indexed sequences; without the choice of $\epsilon$ it is \lean{HPC.Fold.B78.fold_low_end}). With $\epsilon$ chosen by Lemma~\ref{lem:first-direction} it is \lean{HPC.Fold.B79.finfold_low_any}.\\
Fold with a fake high end (Lemma~\ref{lem:fake-high}) & \lean{HPC.Fold.B79.finfold_high}, from \lean{HPC.Fold.B79.fold_high_end_mode}.\\
Lemma~\ref{lem:meet-low} & \lean{HPC.Fold.B79.meet_low_gen}, \lean{HPC.Fold.B79.maxIdx_near_zero}; the two sequences and their properties are collected in \lean{HPC.Fold.B79.seq_facts}.\\
Exclusion (Lemma~\ref{lem:exclusion}) & \lean{HPC.Fold.B79.pair_exclusion}.\\
Prefix fold (Proposition~\ref{prop:prefix-fold}) & \lean{HPC.Fold.B79.pair_prefix_fold}; the conclusion is \lean{HPC.Fold.B79.PrefixFold}.\\
Suffix fold (Corollary~\ref{cor:suffix-fold}) & \lean{HPC.Fold.B79.pair_suffix_fold}.\\
Splice (Proposition~\ref{prop:splice}) & \lean{HPC.Fold.B79.splice_schedule_pair} (the version with $\lambda=\mu=\lo$ and without (c), (d) is \lean{HPC.Fold.B78.splice_schedule}); the passage from levels to links is \lean{HPC.Fold.B78.link_sub_of_level}.\\
From a schedule to the block map (Lemma~\ref{lem:assemble}) & \lean{HPC.Fold.B79.assemble_pair}.\\
Two windows (Lemma~\ref{lem:windows}) & \lean{HPC.Fold.B79.pair_windows}.\\[4pt]
\multicolumn{2}{@{}l}{\emph{Chapter~\ref{ch:descriptions}}}\\
Bing 1951, Theorem~1, body of the proof (Theorem~\ref{thm:desc-exists}) & (a) is \lean{HPC.Plane.PA51.exists_description_sub}; (b) is \lean{HPC.Plane.exists_description_of_chainable_HI} (in \lean{HPC/Plane/PseudoArc51.lean}), via \lean{HPC.Plane.PA51.lift_description}.\\
Bing 1951 (Theorem~\ref{thm:eventually-crooked}) & \lean{HPC.Plane.PA51.eventually_crookedIn} (in \lean{HPC/Plane/PA51Crook.lean}).\\
Lemma~\ref{lem:fail-data} & \lean{HPC.Plane.PA51.FailData} (which records the property (d) through small sets, the links, none inside $C_{k-1}$, resp.\ $C_{h+1}$), \lean{HPC.Plane.PA51.failData_of_fail_lt} (the case $i<j$) and \lean{HPC.Plane.PA51.failData_of_not_crooked}.\\
Lemma~\ref{lem:limit-tools} & \lean{HPC.Plane.PA51.chainedAt_union_singleton}, \lean{HPC.Plane.PA51.limSup_miss} (with \lean{HPC.Plane.PA51.approx_of_mem_limSup}), \lean{HPC.Plane.PA51.exists_limSup_near}.\\
Lemma~\ref{lem:side} & \lean{HPC.Plane.PA51.subset_side_of_miss}, with \lean{HPC.Plane.PA51.lowU} and \lean{HPC.Plane.PA51.highU}.\\
Proposition~\ref{prop:limit-false} & \lean{HPC.Plane.PA51.limit_false}.\\
Definition~\ref{def:predesc} & \lean{HPC.Plane.PA51.PreDesc}.\\
Lemma~\ref{lem:crooked-next} & \lean{HPC.Plane.PA51.exists_crooked_next}.\\
Proposition~\ref{prop:predesc-exists} & \lean{HPC.Plane.PA51.exists_preDesc}.\\
Lemma~\ref{lem:predesc-basic} & \lean{HPC.Plane.PA51.PreDesc.refines_of_le}, \lean{HPC.Plane.PA51.PreDesc.crooked_of_lt}, \lean{HPC.Plane.PA51.PreDesc.eventually_five_le}.\\
Reaching both ends (Lemma~\ref{lem:meets-ends}) & \lean{HPC.Plane.PA51.PreDesc.exists_subchain_meets_ends}, with the limit step \lean{HPC.Plane.PA51.PreDesc.false_of_disjoint_limit}.\\
One step (Lemma~\ref{lem:desc-step}) & \lean{HPC.Plane.PA51.PreDesc.exists_step}; the analogous step for descriptions is \lean{HPC.Description.exists_step} (in \lean{HPC/PseudoArc/T13Step.lean}).\\
Theorem~\ref{thm:desc-univ} & \lean{HPC.Plane.PA51.PreDesc.exists_description}, and with Proposition~\ref{prop:predesc-exists}, \lean{HPC.Plane.PA51.exists_description_univ}.\\
Lemma~\ref{lem:ext} & \lean{HPC.Plane.PA51.ext}, \lean{HPC.Plane.PA51.mem_ext_iff}, \lean{HPC.Plane.PA51.ext_disjoint}, \lean{HPC.Plane.PA51.ext_mono}, \lean{HPC.Plane.PA51.ext_subset_thickening}, \lean{HPC.Plane.PA51.diam_ext_le}, \lean{HPC.Plane.PA51.exists_closure_ext_subset}.\\[4pt]
\multicolumn{2}{@{}l}{\emph{Chapter~\ref{ch:main}}}\\
Pseudo-arc (\vone{Definition~2.14}) & \lean{HPC.Plane.IsPseudoArc}: the conjunction of \lean{HPC.Plane.IsSnakeLike} (which contains ``$M$ is a continuum''), \lean{HPC.Plane.IsHereditarilyIndecomposable} and \lean{Set.Nontrivial}.\\
Bing 1951, Theorem~1, the ``in fact'' clause (Theorem~\ref{thm:unique-points}) & \lean{HPC.Plane.theorem1_HIC}. Formally the described sets are subsets $\Delta_1.M\subseteq M$ and $\Delta_2.M\subseteq M'$ of the spaces $M$, $M'$, equal to the whole spaces; the homeomorphism of \lean{HPC.theorem12} is composed with the identifications $M\cong\Delta_1.M$ and $\Delta_2.M\cong M'$.\\
Uniqueness of the pseudo-arc (Theorem~\ref{thm:unique}) & \lean{HPC.Plane.theorem1_HIC'}; the same statement with the three defining properties written out is \lean{HPC.Plane.theorem1_snake} (Bing, \emph{Snake-like continua}, Theorem~1).\\
Homogeneity of the pseudo-arc (Theorem~\ref{thm:homog}) & \lean{HPC.Plane.homogeneous_of_isPseudoArc}.\\
Corollary~\ref{cor:plane-unique} & \lean{HPC.Plane.isPseudoArc_plane} and \lean{HPC.Plane.pseudoArc_unique}.\\
Lemma~\ref{lem:sub-pseudoarc} & \lean{HPC.Plane.isPseudoArc_of_subcontinuum}.\\
Corollary~\ref{cor:sub-homeo} & \lean{HPC.Plane.homeomorphic_subcontinuum}.\\
Definition~\ref{def:endpoint} & \lean{HPC.Plane.IsEndPoint}, \lean{HPC.Plane.OppositeEndPoints} (in \lean{HPC/Plane/EndPoints51.lean}).\\
Proposition~\ref{prop:opposite} & \lean{HPC.Plane.oppositeEndPoints_of_composants}.\\
Lemma~\ref{lem:endpoint-nested} & \lean{HPC.Plane.condB_of_endPoint}.\\
Lemma~\ref{lem:not-HI-not-end} & \lean{HPC.Plane.exists_not_endPoint_of_not_HI}.\\
Bing 1951, \emph{Snake-like continua}, Theorem~16 (Theorem~\ref{thm:endpoints}) & \lean{HPC.Plane.isPseudoArc_iff_endPoints}.\\
\end{longtable}
